% Generated draft: requires mathematical and source review.
\documentclass[11pt]{article}
\usepackage[margin=25mm]{geometry}
\usepackage{fontspec}
\setmainfont{Latin Modern Roman}
\usepackage{amsmath,amssymb,mathtools,unicode-math}
\setmathfont{Latin Modern Math}
\usepackage{enumitem,needspace,xurl,hyperref}
\setlength{\parindent}{0pt}
\setlength{\parskip}{5pt}
\setlength{\emergencystretch}{3em}
\newcommand{\N}{\mathbb{N}}
\newcommand{\Z}{\mathbb{Z}}
\newcommand{\Q}{\mathbb{Q}}
\newcommand{\R}{\mathbb{R}}
\newcommand{\C}{\mathbb{C}}
\newcommand{\F}{\mathbb{F}}
\newcommand{\A}{\mathbb{A}}
\newcommand{\PP}{\mathbb{P}}
\newcommand{\E}{\mathbb{E}}
\newcommand{\eps}{\varepsilon}
\newcommand{\dd}{\,\mathrm{d}}
\begin{document}
\section*{UnsolvedMath ID: 30005143}
\textbf{Source problem code:} catalogue entry
\subsection{Definitions and mathematical statement}
A Riemannian metric on a smooth manifold is homogeneous if its isometry group acts transitively. A manifold is contractible if its identity map is homotopic to a constant map. Here Ricci flow is interpreted as the unnormalized equation
\[
\frac{\partial g(t)}{\partial t}=-2\operatorname{Ric}(g(t)).
\]
For a homogeneous initial metric \(g_0\), let \(g(t)\), \(0\leq t<T_{\max}\), denote its maximal forward homogeneous Ricci-flow solution with \(g(0)=g_0\). Homogeneity is understood in the usual invariant sense: a transitive group of isometries of \(g_0\) preserves the evolving metrics. The transitive group is not prescribed uniformly across all initial metrics.

\textbf{Dynamical Alekseevskii conjecture (supplied formulation).} For every finite-dimensional, simply connected homogeneous manifold \(M^n\) that is not contractible, and every homogeneous Riemannian initial metric \(g_0\) on \(M\), the maximal forward homogeneous Ricci flow satisfies
\[
T_{\max}(M,g_0)<\infty.
\]
Both compact and noncompact manifolds are included. No positivity assumption is imposed on the scalar curvature or Ricci tensor of \(g_0\).

The supplied source uses ``finite extinction time'' without defining it. The statement above interprets this as finite maximal forward existence time for the unnormalized flow; it does not assert convergence of the entire manifold to a point. The flow normalization and invariant interpretation of homogeneity should be checked against the complete report.

As background, the supplied excerpt asserts that every such \(M\) admits a homogeneous metric \(g_+\) of positive scalar curvature, for each homogeneous presentation \(M=G/H\). This is recorded as a source assertion, not proved here. The conjecture concerns every other homogeneous initial metric as well, including metrics of negative scalar curvature or negative definite Ricci tensor when such metrics exist.
\subsection{Short English statement}
The conjecture asks whether every homogeneous Ricci flow on a simply connected, noncontractible homogeneous manifold has finite maximal forward existence time. The initial metric need not have positive curvature. Here the flow is interpreted as unnormalized; the source's precise meaning of ``extinction'' requires review.
\subsection{Sources}
\begin{itemize}
\item Local UnsolvedMath registry, supplied complete record: permanent ID \(30005143\), exact source problem code \texttt{catalogue entry}, title ``Dynamical Alekseevskii Conjecture for Homogeneous Ricci Flows.'' The registry supplies the status \texttt{open} and statement status \texttt{corrected\_verified}; these labels are not independent mathematical verification. Its statement-verification note says that only formula delimiters, notation, grammar, and typography were repaired, preserving mathematical content and scope.
\item Supplied report citation: \emph{Geometrie (2023)}, Oberwolfach Reports, pp.\ 1551--1601, DOI \texttt{10.4171/owr/2022/28}, \url{https://doi.org/10.4171/owr/2022/28}. The registry separately records proposed year 2022. The excerpt associates the discussion with ``Non-compact Einstein-manifolds with symmetry,'' Christoph B\"ohm, joint work with Ramiro Lafuente; no separate proposer is specified. The excerpt ends with an unfinished ``Clearly,'' followed by contribution information.
\item Supplied literature source, ``OWR report'': \url{https://ems.press/content/serial-article-files/46965}.
\item Supplied literature source, ``Araujo, finite extinction for a family'': \url{https://arxiv.org/abs/2507.05097}.
\item Supplied literature source, ``Billon, dimension five'': \url{https://arxiv.org/abs/2507.17781}.
\item Literature checked at: \texttt{2026-08-22}. This is supplied dated curation, not independently verified. Its assessment labels the conjecture open in full generality, reports finite-extinction results for compact spaces and broad noncompact families, and reports a separate 2025 proof in dimension five. The supplied background explicitly requests independent verification of the time-sensitive status and cited claims; no such verification is claimed here.
\end{itemize}
\end{document}
